\chapter{Ureteral Stent Placement}

\section*{Introduction \& Clinical Indications}
A ureteral stent is a hollow tube placed between the renal pelvis and bladder to maintain urine drainage. It is used for obstruction from stones, strictures, tumors, edema, or postoperative injury; to protect drainage after selected ureteral procedures; or to facilitate staged treatment. In an infected obstructed system, urgent drainage by ureteral stent or nephrostomy and antibiotics takes priority over definitive stone removal.

Stents are usually temporary. The planned removal or exchange date must be documented because forgotten stents can encrust, migrate, fracture, obstruct, and cause renal injury. Stent selection depends on ureteral length, obstruction, expected duration, infection, and whether an external string is appropriate.

\section*{Relevant Surgical Anatomy}
The ureter crosses the pelvic brim near the iliac vessels and enters the bladder obliquely through the intramural tunnel. Physiologic narrowing occurs at the ureteropelvic junction, pelvic brim, and ureterovesical junction. The bladder trigone, ureteral orifices, pelvic organs, and nearby vessels are relevant during cystoscopic and fluoroscopic placement.

\section*{Preoperative Workup \& Preparation}
Review symptoms, renal function, urinalysis and culture when indicated, imaging, allergies, anticoagulants, prior ureteral surgery, and the likely cause and level of obstruction. Treat sepsis urgently and avoid elective manipulation of an infected obstructed system. Explain hematuria, frequency, urgency, flank discomfort, migration, infection, perforation, need for repeat procedures, and the necessity of timely stent removal or exchange.

\section*{Surgical Technique \& Procedure}
Under appropriate anesthesia, cystoscopy identifies the ureteral orifice. A guidewire is advanced under direct and, when needed, fluoroscopic guidance. After the wire reaches the renal pelvis, a stent is advanced so its proximal and distal curls sit in the kidney and bladder. Difficult anatomy may require retrograde contrast, dilation, ureteroscopy, or an alternative nephrostomy. Drainage and position are confirmed, and the stent and follow-up plan are recorded.

\section*{Complications \& Risks}
Early complications include bleeding, infection, ureteral perforation, false passage, stent malposition, migration, and failure to drain. Common stent symptoms include urinary frequency, urgency, dysuria, hematuria, and flank pain during voiding. Prolonged dwell time can cause encrustation, stone formation, fracture, obstruction, and loss of renal function.

\section*{Postoperative Care \& Recovery}
Provide symptom-control advice and clear emergency instructions for fever, rigors, severe pain, vomiting, inability to urinate, or reduced urine output. Confirm that obstruction and infection are improving. Remove or exchange the stent at the planned interval, and ensure the patient understands who is responsible for tracking it. Definitive treatment of the underlying stone, stricture, or tumor should not be lost during temporary drainage.

\section*{Outcomes \& Prognosis}
Stenting is usually effective for restoring drainage and buying time for definitive management. Prognosis depends on the cause and duration of obstruction, infection, renal reserve, and reliable follow-up. A documented stent registry or reminder system helps prevent retained devices.

\section*{References}
\begin{enumerate}
\item European Association of Urology. Guidelines on Urolithiasis.
\item American Urological Association. Surgical Management of Stones: Clinical Guidelines.
\item National Institute for Health and Care Excellence. Renal and ureteric stones: assessment and management.
\end{enumerate}
